\documentclass[10pt]{article}
\usepackage{icmlko}

\icmlmeta{IP 파운데이션 월드모델}{478}{자를 자로 쟀다}{53노트 묵은 채택 문턱을 널 팔 쌍으로 처음 측정했다 --- 옛 자는 헐거웠고, 어제 세운 예리한 자는 자격이 없었다}

\begin{document}

\twocolumn[
\icmlheader

\begin{abstract}
이 실험실은 백 노트 넘게 하나의 상수로 채택을 결정해 왔다 --- 판 $\Delta$ 가
$0.0045$ 를 넘으면 채택. 그런데 그 수는 \textbf{열한 도메인 시대(유보 3,369)의
값}이고 현재 분모는 \textbf{3,775(열두 도메인)}다. 53노트 전에 *"짝
부트스트랩으로 다시 재라"* 고 등록해 두고 한 번도 이행하지 않았다.

직전 사이클이 그 미이행의 대가를 실측했다 --- 같은 결론이 자에 따라
\textbf{36\% 도 되고 146\% 도 된다}. 아이돌 팔의 행 부트 상한이 요구치의
몇 퍼센트인가가, 어느 자를 쓰느냐로 뒤집힌다.

이번 사이클은 문턱의 \textbf{정의부터} 다시 세웠다. \emph{채택 문턱 = 두 팔이
실제로 같을 때 판 $\Delta$ 가 흔들리는 폭의 $2\sigma$}이며, 이 정의는 씨앗
잡음과 표본 잡음을 \textbf{둘 다} 담아야 자격이 있다 --- 하나만 담은 자는
\textbf{실격}으로 측정 전에 못박았다. 그리고 참 $\Delta{=}0$ 인 널 팔 쌍을
만들었다: 같은 팔의 씨앗 0--5 앙상블 대 6--11 앙상블.

결과는 \textbf{R5 $= 0.00353$}이다. 옛 상수의 \textbf{0.784배}로 \emph{옛 자는
헐거웠고}, 어제 세운 예리한 자($0.0011$)의 \textbf{3.21배}로 \emph{그것은
재현성 자이지 채택 문턱이 아니었다}. 새 자로 재면 36\% 대 146\% 는
\textbf{45.5\%} 로 결착된다 --- 여전히 못 넘는다. 그리고 열두 도메인 중
\textbf{쓸모 문턱을 넘는 것은 어느 자로도 0개}다.
\end{abstract}
]

\section{53노트 동안 안 잰 것}

이 판의 채택 규칙은 한 줄이다 --- \emph{판 $\Delta \ge 0.0045$ 면 채택}.
그런데 그 상수의 출처는 \textbf{유보 3,369행(열한 도메인)} 이고, 열두 번째
도메인이 합류한 뒤 분모는 \textbf{3,775} 다. 미결로 등재된 지 53노트가 됐다.

직전 사이클이 그 방치의 비용을 처음 값으로 보여 줬다. 아이돌 팔의 행 부트
상한 $0.11877$ 이 \emph{요구 도메인 $\Delta$} 의 몇 퍼센트인가는:

\begin{center}\small
\begin{tabular}{lrr}
\toprule
자 & 요구 $\Delta$ & 상한의 몫 \\
\midrule
옛 상수 $0.0045$ & 0.3331 & 35.7\% \\
직전 사이클의 예리한 자 $0.0011$ & 0.0814 & \textbf{145.9\%} \\
\bottomrule
\end{tabular}
\end{center}

\textbf{한쪽은 ``가망 없다'' 이고 다른 쪽은 ``넘는다'' 다.} 같은 자료·같은 팔인데
자만 바꿨다.

\claimx{문턱을 안 재고 두면, 문턱이 결론의 부호를 정한다.}{새 자로 쟀을 때
두 값 사이가 아니라 어느 한쪽에 딱 붙으면, 이 걱정이 과했던 것이다.}

\section{정의부터 다시}

측정 전에 정의를 못박았다.

\begin{quote}
\textbf{채택 문턱}은 \emph{두 팔이 실제로 같을 때} 판 $\Delta$ 가 흔들리는
폭의 $2\sigma$ 다. 이 정의는 \textbf{씨앗 잡음과 표본 잡음을 둘 다} 담아야
자격이 있다.
\end{quote}

그리고 \textbf{실격 규칙}을 같이 적었다 --- 씨앗만 담은 자도, 행만 담은 자도
\textbf{단독으로는 채택 문턱 자격이 없다}. 이 규칙이 직전 사이클이 세운
``예리한 자''($0.0011$, 12씨앗 짝 평균의 $2\sigma$)를 \textbf{스스로 실격}시킨다.
자기가 만든 자를 측정 전에 배제하는 것이 요점이었다.

\textbf{널 팔 쌍}이 이 정의를 실현한다. 같은 팔($K{=}1$)의 씨앗 0--5 순위평균
앙상블과 6--11 앙상블을 두 팔로 놓으면 \textbf{참 $\Delta$ 가 0} 이다. 그
쌍에 도메인 층화 행 부트($B{=}10{,}000$)를 걸어 흔들림을 잰다. $K{=}2$ 에서도
같이 재고 \textbf{큰 쪽을 채택}한다 --- 채택은 비가역이므로 보수적으로 간다.

\textbf{배선을 먼저 닫았다.} 직전 사이클 산출물을 sha256 으로 고정하고
(손 전사 금지), 12씨앗 $\times$ 2팔 $\times$ 12도메인 = \textbf{288칸을
재계산해 불일치 0}, 그리고 그 사이클의 행수준 부트를 \textbf{소수 다섯 자리까지
재현}했다($-0.00013$, $[-0.00326, +0.00306]$, SD $0.00161$).

\section{자 다섯}

\begin{figure}[t]
\centering
\includegraphics[width=\columnwidth]{figs/fig1.pdf}
\caption{왼쪽: 자 다섯의 눈금(로그). 오른쪽: 같은 결론(아이돌 행 부트 상한이
요구치의 몇 \%)이 자에 따라 어디까지 움직이나.}
\label{fig:1}
\end{figure}

\begin{center}\small
\begin{tabular}{lrl}
\toprule
자 & $2\sigma$ & 성질 \\
\midrule
R2 씨앗 짝(n=12) & 0.00110 & 씨앗만 --- \textbf{실격} \\
\textbf{R5 합성} & \textbf{0.00353} & \textbf{씨앗 $\oplus$ 행 짝} \\
R1 씨앗 수준(단일) & 0.00408 & 씨앗만 --- 실격 \\
옛 상수 & 0.00450 & 11도메인 유물 \\
R3 행 수준(비짝) & 0.02628 & 행만 --- 실격 \\
\bottomrule
\end{tabular}
\end{center}

R5 의 성분은 씨앗 $0.00099$(12대12 환산)와 행 짝 $0.00146$ 이다. 널 쌍의
$2\sigma$ 는 $K{=}1$ 에서 $0.00288$, $K{=}2$ 에서 $0.00292$ 로 거의 같았고
보수적으로 후자를 썼다.

\claimx{채택 문턱은 $0.00353$ 이다. 옛 상수는 그 1.27배로 헐거웠고, 직전
사이클의 예리한 자는 그 0.31배로 채택 판정에 쓸 자격이 없다.}{널 쌍의
$2\sigma$ 가 팔에 따라 1.5배 넘게 벌어지면 ``팔에 무관한 상수 문턱'' 자체가
성립하지 않는다 --- 실측은 $0.00288$ 대 $0.00292$ 로 1.01배였다.}

\section{그래서 36\% 대 146\% 는}

\begin{center}\small
\begin{tabular}{lrr}
\toprule
자 & 요구 아이돌 $\Delta$ & 상한의 몫 \\
\midrule
옛 $0.0045$ & 0.3331 & 35.7\% \\
\textbf{R5 $0.00353$} & \textbf{0.2611} & \textbf{45.5\%} \\
예리한 $0.0011$ & 0.0814 & 145.9\% \\
\bottomrule
\end{tabular}
\end{center}

\textbf{옳게 잰 자로는 45.5\% 이고 여전히 못 넘는다.} 직전 사이클이 만든
$146\%$ 쪽이 틀린 쪽이었다 --- 그 자가 표본 잡음을 하나도 안 담았기 때문이다.
사전등록의 예측(32--62\%)이 적중했다.

\section{열두 도메인은 어떻게 되나}

문턱을 두 칸으로 나눠 채점했다 --- \textbf{칸1 ``0 이 아닌가''}
($|\Delta|/\text{짝SD} \ge 1.0$ \emph{그리고} $|\Delta|/\text{짝SE} \ge 3.60$,
후자는 12검정 본페로니), \textbf{칸2 ``쓸모 있나''}
($|\Delta_d \cdot w_d| / \sum w \ge$ 문턱).

\begin{center}\small
\begin{tabular}{lrrc}
\toprule
 & $|\Delta|/$SE & 판 기여 & 칸1 \\
\midrule
게임 & 4.56 & 0.000246 & \checkmark \\
애니 & 4.37 & 0.000371 & \checkmark \\
펀딩 & 3.89 & 0.000551 & \checkmark \\
아이돌 & 3.70 & 0.000361 & \checkmark \\
도서 & 3.49 & 0.000649 & \\
팝업 & 2.78 & 0.000243 & \\
영화 & 2.65 & 0.000277 & \\
세계애니 & 2.58 & 0.000244 & \\
만화 & 2.28 & 0.000133 & \\
웹툰 & 2.04 & 0.000641 & \\
모바일 & 1.11 & 0.000068 & \\
시장팝업 & 0.85 & 0.000044 & \\
\bottomrule
\end{tabular}
\end{center}

\textbf{칸1 은 4/12, 칸2 는 0/12 다.} 그리고 칸2 는 \textbf{어느 자로 재도
0/12} 다 --- R5 로도, 옛 $0.0045$ 로도, 예리한 $0.0011$ 로도. 부족 배수는
웹툰 5.5배부터 시장팝업 79.5배까지다.

직전 사이클이 자백한 ``$|\Delta|/\text{SE} \ge 2$ 가 10/12 를 통과시킨다'' 는
결함은 칸1 을 **4/12** 로 좁혀 고쳐졌다. 그러나 그것은 \emph{0 이 아니다}를
물을 뿐이고, \emph{쓸모 있나}는 여전히 아무도 통과하지 못한다.

\claimx{특기 칸을 하나 더 주는 팔은 어느 자로도 채택 문턱을 못 넘는다
($|\Delta|/\text{R5} = 0.102$).}{유보 행이 크게 늘거나 부호가 한 방향으로
정렬되면 같은 팔이 칸2 에 접근해야 한다.}

\section{남는 것}

이 사이클도 채택을 하나도 안 냈다. 남긴 것은 \textbf{자} 하나다 ---
$0.0045$ 를 $0.00353$ 으로 바꿨고, 그 수가 **어디서 왔는지**가 처음으로
문서에 남았다(널 팔 쌍 · 씨앗 $\oplus$ 행 · 보수적 팔 선택).

그리고 이 사이클이 확인한 것 하나가 더 있다. 직전 사이클은 자기가 만든
``예리한 자''를 다음 사이클의 후보로 남겼는데, \textbf{그 자는 실격이었다}.
결과를 본 뒤에 자를 고르면 자가 결론을 따라간다 --- 그래서 실격 규칙을
\textbf{측정 전에} 적었고, 그것이 이번 사이클에서 실제로 발동했다.

\end{document}
